% Generated from locale/tr/content/model-theory/basics/theory-of-m.tex; do not hand-edit.


\olfileid[tr]{mod}{bas}{thm}

\section{Bir \printtoken{S}{structure} Kuramı}

Her !!{structure}~$\Struct M$ bazı !!{sentence}s ifadelerini doğru,
bazılarını yanlış kılar. Doğru kıldığı bütün !!{sentence}s kümesine
onun \emph{kuramı} denir. Bu küme gerçekten de bir kuramdır; çünkü
mantıksal olarak gerektirdiği her şey, $\Struct M$ dâhil bütün
modellerinde doğru olmak zorundadır.

\begin{defn}
  !!a{structure}~$\Struct M$ verildiğinde $\Struct M$ yapısının
  \emph{kuramı}, $\Struct M$ içinde doğru olan !!{sentence}s
  kümesi~$\Theory M$'dir; yani
  $\Theory M=\Setabs{!A}{\Sat{M}{!A}}$.
\end{defn}

``Kuram'' sözcüğünü, türetim bakımından kapalı olsun ya da olmasın,
amaçlanan bir yorumu bulunan !!{sentence}s kümeleri için de gayriresmî
olarak kullanırız.

\begin{prop}
Her $\Struct M$ için $\Theory M$ tamdır.
\end{prop}

\begin{proof}
Her !!{sentence}~$!A$ için ya $\Sat{M}{!A}$ ya da
$\Sat{M}{\lnot !A}$; dolayısıyla ya $!A\in\Theory M$ ya da
$\lnot !A\in\Theory M$.
\end{proof}

\begin{prop}\ollabel{prop:equiv}
  $\Theory M$ içindeki her $!A$ için $\Struct N\models !A$ ise
  $\Struct M\elemequiv\Struct N$.
\end{prop}

\begin{proof}
$\Theory M$ içindeki bütün $!A$ için $\Sat{N}{!A}$ olduğundan
$\Theory M\subseteq\Theory N$. $\Sat{N}{!A}$ ise
$\Sat/{N}{\lnot !A}$, dolayısıyla $\lnot !A\notin\Theory M$.
$\Theory M$ tam olduğundan $!A\in\Theory M$. Böylece
$\Theory N\subseteq\Theory M$ ve $\Struct M\elemequiv\Struct N$.
\end{proof}

\begin{rem}\ollabel{remark:R}
  Evreni gerçel sayılar kümesi~$\Real$ olan ve yalnızca gerçeller
  üzerindeki $<$ bağıntısı olarak yorumlanan iki yerli bir
  !!{predicate} içeren !!{language} içindeki
  $\Struct R=\langle\Real,<\rangle$ !!{structure} ifadesini ele alalım.
  $\Struct R$ açıkça !!{nonenumerable}; fakat $\Theory R$ elbette
  tutarlı olduğundan Löwenheim--Skolem teoremi gereği, örneğin
  $\Struct S$ diyeceğimiz !!a{enumerable} modele sahiptir ve
  \olref{prop:equiv} gereği $\Struct R\equiv\Struct S$. Üstelik
  $\Struct R$ ile $\Struct S$ izomorf olmadığından bu,
  \olref[iso]{thm:isom} teoreminin tersinin genel olarak geçersiz
  olduğunu gösterir.
\end{rem}

